% !Mode:: "TeX:UTF-8" 

\BiChapter{绪论}{Introduction}

\BiSection{研究背景及意义}{Research Background and Significance}

随着信息技术的飞速发展，云计算已经成为各行各业数字化转型的关键驱动力。其弹性扩展、按需付费以及便捷高效的特性，使得越来越多的组织选择将数据、应用程序和服务迁移至云端\cite{IAM_in_Cloud}。云计算的普及极大地提升了业务的灵活性和可扩展性，但也随之带来了前所未有的安全挑战。由于数据存储在远程的云基础设施中，并且可以通过互联网随时随地访问，传统的安全边界变得模糊，使得保护云环境下的敏感信息和资源变得至关重要\cite{Why_cloude_security_is_important}。

在众多的云安全问题中，身份与访问管理 (Identity and Access Management, IAM) 占据着核心地位。作为云服务提供商提供的一项基础安全服务，IAM 负责管理用户和应用程序对云资源和服务的访问权限，确保只有经过授权的实体才能在适当的时间以适当的方式访问相应的资源\cite{googleIdentityAccess}。如果 IAM 系统设计不当或配置错误，可能会导致未经授权的访问，进而引发数据泄露、服务中断等严重的后果\cite{Why_need_IAM}。因此，深入研究云计算环境下的 IAM 技术和安全控制机制，对于保障云服务的安全稳定运行具有重要的现实意义。

云计算环境下的 IAM 与传统的身份管理有所不同，它需要应对更加复杂和动态的场景。这包括管理分布在不同地理位置的资源、处理各种类型的用户（例如员工、客户、合作伙伴）以及非人类实体（例如应用程序、服务）的访问需求\cite{Model_IAM}。此外，随着网络安全威胁的日益演变和数据保护法规的日益严格，传统的基于密码的静态 IAM 已经难以满足现代云安全的要求，需要采用更加动态和自适应的方法\cite{digitaloceanWhatCloud}。例如，频繁发生的数据泄露事件和账户劫持事件表明，仅仅依赖用户名和密码进行身份验证是远远不够的\cite{Why_cloude_security_is_important}。同时，诸如 通用数据保护条例 (General Data Protection Regulation, GDPR)、健康保险便携与责任法案 (Health Insurance Portability and Accountability Act, HIPAA) 和支付卡行业数据安全标准 (Payment Card Industry Data Security Standard, PCI DSS) 等法规对云端数据的保护提出了明确的要求，也进一步凸显了强化云 IAM 的必要性\cite{Why_cloude_security_is_important}。因此，对云计算环境下 IAM 的研究不仅具有重要的学术价值，也对指导实际应用和提升云安全水平具有深远的意义。

% \BiSection{身份访问管理的作用}{The Role of IAM}

\BiSubsection{云计算模型及其安全影响概述}{}

云计算提供了多种服务模型，每种模型在责任划分和安全考量上都有所不同\cite{IAM_Concepts}。基础设施即服务 (IaaS) 提供最基础的计算资源，例如虚拟机和存储，用户对操作系统、应用程序和数据拥有最大的控制权，因此也承担了大部分的安全责任\cite{IAM_Concepts}。在 IaaS 模式下，有效的 IAM 配置对于控制对这些关键组件的访问至关重要。平台即服务 (PaaS) 在 IaaS 的基础上提供了开发和部署应用程序的平台，云服务提供商负责底层基础设施的安全，而用户则需要关注其应用程序的安全以及对应用程序和平台服务的访问控制\cite{IAM_Concepts}。软件即服务 (SaaS) 直接通过互联网提供应用程序，云服务提供商承担了大部分的安全责任，但用户仍然需要管理其在应用程序内的用户账户和权限\cite{IAM_Concepts}。

除了服务模型，云计算还存在不同的部署模型，包括公有云、私有云、混合云和联邦云\cite{IAM_Concepts}。公有云在多租户环境下运行，需要强大的 IAM 控制来隔离不同用户的数据和资源。私有云提供更高的控制权，但仍然需要有效的 IAM 来管理组织内部的访问。混合云和联邦云涉及跨多个环境管理身份和访问，带来了额外的复杂性，通常需要采用身份联邦等技术来实现统一管理\cite{IAM_Concepts}。不同的云计算模型在安全责任的分配上存在显著差异，这直接影响了需要实施的 IAM 策略\cite{IAM_Concepts}。理解这些责任边界对于设计和实施有效的云 IAM 框架至关重要。

\BiSubsection{云计算环境中的主要安全挑战}{}

云计算环境面临着诸多安全挑战，其中数据泄露是首要关注的问题\cite{Why_cloude_security_is_important}。由于大量敏感信息存储在云端，一旦发生未经授权的访问或数据被窃取，将会对组织造成严重的损失\cite{Why_cloude_security_is_important}。弱口令、权限配置不当等 IAM 方面的疏忽往往是导致数据泄露的重要原因\cite{Use_Secure_IAM}。云配置错误也是一个常见的安全风险\cite{Why_cloude_security_is_important}。例如，IAM 相关的配置错误可能导致不应被公开的资源暴露在互联网上，为攻击者提供了可乘之机\cite{Why_cloude_security_is_important}。未经授权的访问是 IAM 需要解决的核心问题\cite{Why_cloude_security_is_important}。IAM 通过验证用户身份和限制其操作权限来防止未经授权的实体访问云资源。

内部威胁也是云安全中不可忽视的风险\cite{bigidEnhanceSecurity}。即使是拥有合法访问权限的内部人员，也可能出于恶意或疏忽而滥用其权限，导致安全事件的发生\cite{bigidEnhanceSecurity}。账户劫持是另一种常见的威胁\cite{bigidEnhanceSecurity}。攻击者通过各种手段（例如网络钓鱼、恶意软件）获取合法用户的账户凭据，然后冒充该用户访问云资源\cite{bigidEnhanceSecurity}。在采用多云策略的组织中，跨多个云环境管理身份变得尤为复杂\cite{NavigatingChallenges}。不同的云服务提供商拥有各自的 IAM 系统，如何确保在这些异构环境中实施一致的安全策略是一个巨大的挑战\cite{NavigatingChallenges}。此外，现代云应用不仅涉及人类用户，还包括大量的非人类实体，例如应用程序、服务和 API，如何有效地管理这些非人类实体的身份和访问权限也变得越来越重要\cite{cloudAllianceIdentityAccess}。随着员工在组织内角色的变化，其访问权限可能会不断累积，导致权限过剩，这被称为角色蔓延\cite{nordlayerBenefitsChallenges}。这种现象增加了潜在的安全风险，因为用户可能拥有不再需要的敏感数据访问权限。

\BiSubsection{身份访问管理 (IAM) 系统的应对}{}

IAM 通过多种机制来应对上述安全挑战。身份验证是 IAM 的第一步，它负责验证尝试访问云资源的用户或实体的身份\cite{IAM_in_Cloud}。这可以通过多种方法实现，例如密码、多因素身份验证 (MFA)、生物识别和无密码身份验证等\cite{IAM_in_Cloud}。授权是 IAM 的下一步，它确定了经过身份验证的用户或实体对特定云资源的访问级别以及允许执行的操作\cite{IAM_in_Cloud}。授权通常基于预定义的访问控制策略进行\cite{IAM_in_Cloud}。这些策略规定了谁可以在什么条件下访问哪些资源。

集中化管理是云 IAM 的一个重要特点\cite{Why_cloude_security_is_important}。通过统一的平台管理用户身份、权限和安全策略，可以简化管理工作并提高安全性\cite{Why_cloude_security_is_important}。监控和审计是 IAM 的另一个关键功能\cite{Why_cloude_security_is_important}。通过跟踪用户活动和访问模式，可以检测到可疑行为并确保合规性\cite{Why_cloude_security_is_important}。

\BiSubsection{身份访问管理 (IAM) 系统的益处}{}

实施强大的云 IAM 系统可以带来多方面的益处。首先，它可以显著增强安全性，降低未经授权的访问、数据泄露和其他安全事件的风险\cite{Why_cloude_security_is_important}。通过严格控制谁可以访问哪些资源，IAM 可以有效地保护云环境中的敏感数据和资产。其次，IAM 有助于满足各种法规要求\cite{Why_cloude_security_is_important}。例如，GDPR、HIPAA 和 PCI DSS 等法规都对数据保护和访问控制提出了明确的要求，而 IAM 提供的控制和审计功能可以帮助组织遵守这些规定。
此外，IAM 还可以提高运营效率\cite{googleIdentityAccess}。通过简化用户访问管理、减少管理开销和提高用户生产力，IAM 可以帮助组织更高效地运营。例如，单点登录 (SSO) 功能允许用户使用一套凭据访问多个云应用程序，从而减少了管理多个密码的需求并提高了用户体验。强大的 IAM 还可以带来成本节约\cite{identitymanagementinstituteCloudSecurity}。通过降低数据泄露和潜在罚款的风险，以及优化资源利用率，IAM 可以帮助组织节省资金。最后，IAM 可以改善治理和可见性\cite{googleIdentityAccess}。通过增强对用户访问和活动的控制，并提供对潜在风险的更好洞察，IAM 可以帮助组织更好地管理其云环境并降低安全风险。

\BiSection{国内外研究应用现状}{}

IAM 的发展历程反映了 IT 基础架构和安全理念的深刻变革。传统的 IT 安全范式在很大程度上依赖于网络层控制，例如防火墙和入侵检测系统，以构建可信的内部网络边界。在这种模式下，IAM 系统通常是本地部署的，与企业的物理网络紧密集成。然而，云计算的兴起彻底改变了这一格局。企业数据和应用逐渐迁移到云端，远程办公和移动访问成为常态，传统的网络边界变得模糊甚至消失。这使得 IAM 的角色愈发重要，它从一个静态的、面向边界的访问控制机制，演变为一种动态的、以身份为中心的、适应性强的安全策略。与本地部署的 IAM 相比，云环境下的 IAM 管理面临着新的挑战：

\begin{enumerate}[wide, label=\arabic*),  labelindent=21pt]
    \item 复杂性增加：云环境通常涉及多种服务、资源和应用，身份和权限关系错综复杂，管理难度远超本地系统。
    \item 分布式特性：身份和资源可能分布在不同的云平台、地理区域甚至法律管辖区，增加了统一管理的难度。
    \item 变化速度加快：云资源的动态伸缩和应用的快速迭代要求IAM系统能够快速响应身份和权限的变更需求。
\end{enumerate}

在云环境中，IAM的责任由云服务提供商 (Cloud Service Providers, CSPs) 和云服务客户 (Cloud Service Customers, CSCs) 共同承担。CSP 负责提供 IAM 工具和平台，并确保其底层基础设施的安全性；而 CSC 则负责正确配置和使用这些 IAM 工具，定义和管理其用户、角色和策略，以保护其在云上的资源和数据。这种共享责任模型虽然带来了灵活性，但也对客户的 IAM 管理能力和安全意识提出了更高要求。如果客户未能妥善履行其 IAM 配置和管理职责，例如授予了过度的权限，那么即使 CSP 提供了安全的 IAM 平台，整体安全也可能受到损害。这凸显了在云时代，理解和有效管理共享责任模型对于确保IAM成功至关重要。

\BiSubsection{国内应用现状}{}

阿里云的资源访问控制服务 (Resource Access Management, RAM) 是其管理用户身份和资源访问权限的核心组件。RAM 允许主账户创建和管理 RAM 用户（代表员工、系统或应用程序），并通过策略授权来控制这些用户对阿里云资源的操作权限，旨在避免共享主账户密钥，实现最小权限分配，从而降低信息安全风险\cite{aliyunRAM}。

RAM 的核心功能包括：

\begin{enumerate}[wide, label=\arabic*),  labelindent=21pt]
    \item 身份管理：创建和管理 RAM 用户、用户组和角色。RAM 用户本身不拥有资源，其创建的资源归属于主账户，相关费用也由主账户统一支付。
    \item 权限策略：支持系统策略（由阿里云管理）和自定义策略（由用户管理），策略语言用于表达细粒度的授权语义，可以精确到 API 操作级别和特定资源 ID，并支持多种限制条件（如来源IP、访问时间、MFA）。
    \item 访问控制：集中控制 RAM 用户的访问密钥和权限，支持 MFA。
    \item 身份联合：虽然 RAM 本身是阿里云内的身份体系，但阿里云也通过 Keycloak 等开源 IAM 解决方案的集成示例，展示了如何实现与外部 IdP 的 SSO、SAML 2.0、OAuth 2.0 和 OIDC 集成。
\end{enumerate}

\BiSubsection{国外应用现状}{}

AWS IAM 是AWS云平台中用于安全控制对AWS服务和资源访问的核心服务。其核心概念包括用户 (IAM Users)、用户组 (Groups)、角色 (Roles) 和策略 (Policies)。策略以 JSON 格式定义，规定了允许或拒绝的操作和资源。AWS 强调区分根用户 (Root User) 和 IAM 用户，推荐使用 IAM 用户进行日常操作，并遵循最小权限原则\cite{awsIAM}。

AWS IAM的关键特性包括：

\begin{enumerate}[wide, label=\arabic*),  labelindent=21pt]
    \item 多因素认证  (MFA)：为用户登录和敏感操作提供额外安全层。
    \item 临时安全凭证：通过 IAM 角色，应用程序和服务可以获取临时凭证访问AWS资源，避免使用长期访问密钥。
    \item 权限边界 (Permissions Boundaries)：用于限制 IAM 实体（用户或角色）可以拥有的最大权限，即使其身份策略允许更多权限，也无法超越权限边界的限制。
    \item 服务控制策略 (Service Control Policies, SCPs)：在 AWS Organizations 中使用，用于在组织或组织单元（OU）级别对账户的权限进行集中控制。
    \item 基于属性的访问控制 (Attribute-Based Access Control, ABAC)：允许使用标签等属性来定义更动态和细粒度的访问策略，例如结合 SAML 会话标签实现 ABAC。
    \item AWS IAM Identity Center (原AWS SSO)：提供统一的身份联合访问入口，方便用户使用企业身份或社交身份访问 AWS 账户和应用。
\end{enumerate}

Microsoft Entra ID（前身为Azure Active Directory）是微软提供的全面的云身份和访问管理服务，是 Microsoft Entra 产品家族的核心组成部分。它为用户、设备、应用和工作负载提供身份验证和授权服务，支持访问云端和本地资源\cite{msEntra}。

Entra ID的核心功能包括：

\begin{enumerate}[wide, label=\arabic*),  labelindent=21pt]
    \item 身份和目录服务：将用户作为目录对象进行管理，支持组管理和 Azure RBAC（基于角色的访问控制）。
    \item 多因素认证 (MFA) 和 条件访问 (Conditional Access)：MFA 提供强大的身份验证，条件访问策略则允许基于用户、位置、设备状态、应用程序、实时风险等多种因素动态评估和强制执行访问控制。
    \item 身份联合：广泛支持 SAML、OAuth 2.0、OIDC 等标准协议，实现与外部 IdP 的联合身份验证和单点登录。
    \item 混合身份管理：通过 Microsoft Entra ID Connect，可以同步本地 Active Directory 身份到云端，实现混合环境下的统一身份管理和 SSO。
    \item 工作负载身份管理：为应用程序和服务提供安全的身份验证和授权机制。
    \item 无密码认证：支持 FIDO2 安全密钥、Windows Hello、身份验证器应用 (OTP) 等多种无密码登录方式。
\end{enumerate}

\BiSubsection{前沿进展}{}

IAM 领域的研究正以前所未有的速度向前推进，多个前沿方向共同塑造着下一代身份与访问管理的蓝图。这些研究不仅致力于解决当前面临的挑战，更着眼于应对未来可能出现的安全威胁和技术变革。

\textbf{零信任架构 (Zero Trust Architecture, ZTA)}：零信任架构已成为现代网络安全的核心理念之一，它彻底颠覆了传统的基于边界的信任模型，强调“永不信任，始终验证”。ZTA的核心原则包括：对每一次访问请求进行持续验证，无论其来源（内部或外部网络）；实施微隔离以限制攻击的横向移动；严格执行最小权限访问；以及基于用户身份、设备状态、数据敏感性和环境上下文进行动态和情境感知的访问控制决策。

\textbf{去中心化身份与可验证凭证 (VC)}：去中心化身份 (Decentralized Identity, DI) 和可验证凭证 (Verifiable Credentials, VC) 代表了身份管理领域的一项重大范式转变，其核心理念是将数字身份的控制权从中心化的身份提供商交还给用户自身，实现用户对其身份数据的高度自主管理。

\textbf{用于IAM的后量子密码 (PQC)}：量子计算的快速发展对当前广泛依赖于大数分解和离散对数等数学难题的公钥密码体系（如RSA、ECC）构成了潜在的颠覆性威胁。由于 IAM 系统中的数字签名、密钥交换、证书体系等核心安全机制大量使用这些密码算法，因此，量子计算的出现对 IAM 的长期安全性提出了严峻挑战。后量子密码 (Post-Quantum Cryptography, PQC) 的研究旨在开发能够抵抗传统计算机和量子计算机攻击的新型密码算法。

\BiSubsection{对比总结}{}

对比全球与中国的 IAM 发展格局，可以发现两者在技术演进、市场驱动因素和监管影响等方面既有趋同性，也存在显著差异。

从核心技术层面来看，中国 IAM 的发展在很大程度上与全球趋势保持同步。无论是身份认证技术（MFA、SSO、生物识别）、授权模型 (RBAC、ABAC)，还是联合身份协议 (SAML、OAuth、OIDC)，这些基础的 IAM 构建模块在中国市场同样得到了广泛应用和研究。中国本土的云服务提供商（如阿里云、腾讯云、华为云）提供的 IAM 产品，在核心功能上与国际主流 CSP (AWS、Azure、Google Cloud) 的 IAM 服务具有可比性，都致力于提供用户管理、策略控制、角色分配、MFA 等基础能力。

然而，在技术应用的侧重点和创新方向上，中外可能存在差异：

\begin{enumerate}[wide, label=\arabic*),  labelindent=21pt]
    \item 本土化创新：中国市场拥有庞大的用户基数和独特的互联网生态（如微信、支付宝等超级应用的主导地位），这为 IAM 的本土化创新提供了土壤。例如，腾讯云 CAM 支持微信扫码作为 MFA 方式，这种与本土流行应用和用户习惯的深度集成，是国际厂商短期内难以复制的优势。
    \item 对新兴技术的采纳节奏：在全球范围内，ZTA、AI/ML 在 IAM 中的应用、去中心化身份等前沿技术是研发热点。中国在这些领域也有积极探索，例如在数字身份框架中讨论区块链和隐私计算的应用，以及本土云厂商在其安全解决方案中开始融入 AI 能力。但这些技术的实际落地规模和成熟度可能与国际领先水平存在一定差距，部分原因可能在于市场需求、投入以及生态系统的成熟度。
    \item 开源技术的应用与贡献：全球IAM领域受益于众多开源项目（如Keycloak 52, Open Policy Agent等）。中国企业和开发者对开源技术的采纳和贡献也在增加，但在 IAM 核心开源项目的领导力和影响力方面，与国际领先者相比仍有提升空间。
\end{enumerate}

\BiSection{本文主要内容}{}

本文成功复现了基于 SMT 理论的分层抽象算法，通过 Java 语言构建了一个可以部署在云服务器上的 IAM 验证服务，并扩展了原本 AWS 的关键字段，以支持不同云服务厂商的用户的验证需求。

具体而言，本文的主要工作包括：

\begin{enumerate}[wide, label=\arabic*),  labelindent=21pt]
    \item 详细调研了 IAM 及形式化 IAM 验证的相关技术，研究对比了各种 IAM 求解方式，最终选定以基于 SMT 理论的形式化求解为主要研究方向
    \item 基于 Z3 求解器，设计了一套将 IAM 策略语言编码为 SMT 表达式的方案，并扩展 Z3 为 Z3 自动机，以支持对包含正则表达式的 IAM 策略进行编码
    \item 在 Z3 自动机的基础上，进一步设计复现了 AWS 的分层抽象 (Stratified Abstraction)\cite{Stratified}，并讨论了该算法存在的问题以及可能的性能改进
    \item 针对 AWS 提供的异步验证 API, 设计并实现了一套轮询监测的实验流程，用以对比本文的实现方案和 AWS 私有实现方案的正确性和性能
\end{enumerate}

\BiSection{本文组织结构}{}

本文主要分为四个部分，第一部分是相关技术介绍，这里将讲解形式化验证技术的理论、SMT 求解器等内容；第二部分是对形式化验证系统的阐述，包括 SMT 编码器和分层抽象算法；第三部分是系统的实现和验证；第四部分是全文的总结与未来研究方向。

具体章节安排如下：

\begin{enumerate}
    \item 绪论。在这一章中，本文将介绍研究背景和研究意义、国内外研究现状、本文主要工作以及全文组织结构。
    \item 相关技术介绍。在这一章中，本文将考察形式化验证系统的相关理论和技术，包括 SMT 理论，SMT 求解器等。然后本文将考察一个典型的 SMT 求解器：Z3 求解器，指出其理论核心、系统架构、运行策略和主要功能。最后，本文将讨论形式化验证的挑战以及若干替代与补充办法。
    \item SMT 编码器设计。在这一章中，本文将讨论基于 Z3 的 IAM SMT 编码器的编码方法，并扩展 Z3 为 Z3 自动机以支持对正则表达式的编码。
    \item 分层抽象算法设计。在这一章中，本文将讲述分层抽象算法的思想、相关性质、实现算法与理论模型，意在提供一个对该算法的较为详细的讨论。
    \item 形式化验证系统的实现。在这一章中，本文将给出 SMT 编码器和分层抽象算法的实现过程，包括系统中的关键组件的实现以及它们之间的交互过程。
    \item 实验设计与结果。在这一章中，本文将对上一章中实现的形式化验证系统进行若干实验测试，并与已有的云验证系统的结果进行对比验证。
    \item 结论与展望。在这一章中，本文将总结全文内容并讨论算法可能的改进点。
\end{enumerate}
